\documentclass[article,shortnames,nojss]{jss}

%% --- packages ---
\usepackage{amsmath,amssymb,amsthm}
\usepackage{booktabs}
\usepackage{graphicx}
\usepackage{lmodern}
\usepackage{microtype}
\usepackage{orcidlink}
\newcommand{\R}{\mathbb{R}}

%% --- overflow tolerance ---
\setlength{\emergencystretch}{4em}
\hbadness=10000
\hfuzz=20pt

%% --- math macros (avoid \E and \R which jss.cls defines) ---
\providecommand{\Rset}{\mathbb{R}}
\renewcommand{\E}{\mathbb{E}}
\newcommand{\indep}{\perp\!\!\!\perp}
\newcommand{\ones}{\mathbf{1}}
\newcommand{\OTIS}{\textsf{OTIS}}
\newcommand{\SIU}{\textsf{SIU}}
\newcommand{\TPS}{\textsf{TPS}}
\newcommand{\CCRSO}{\textsf{CCRSO}}
\newcommand{\IAP}{\textsf{IAP}}

%% --- bibliography (jss.cls already configures natbib + jss.bst) ---

%% --- title / author block (jss conventions) ---
\title{MRM: Multilevel Reconciliation Methodology \\ \large A multi-source statistical foundation for Canadian carceral, police, and oversight data}
\Plaintitle{MRM: Multilevel Reconciliation Methodology -- A multi-source statistical foundation for Canadian carceral, police, and oversight data}
\Shorttitle{MRM: Multilevel Reconciliation Methodology}

\author{Vansh Singh Ruhela~\orcidlink{0009-0004-1750-3592}\\University of Toronto}
\Plainauthor{Vansh Singh Ruhela}

\Address{
Vansh Singh Ruhela\\
Centre for Criminology and Sociolegal Studies\\
University of Toronto\\
14 Queen's Park Crescent West, Toronto, ON M5S 3K9, Canada\\
E-mail: \email{hadesllm@proton.me}\\
ORCID: \orcidlink{0009-0004-1750-3592}\\
Pronunciation: vansh sing roo-HAY-lah
}

\Abstract{
The MRM (Multilevel Reconciliation Methodology) framework provides a unified
statistical foundation for analysing Canadian carceral, police and
oversight data, integrating five public sources (provincial
\OTIS{} restrictive-confinement microdata, federal $\SIU$
implementation reports, $\IAP$ academic replications, Ontario $\SIU$
police-oversight reports, and the \TPS{} open-data crime corpus).
The framework couples a ten-estimator per-individual causal
ensemble (IPW, AIPW, IRM-DML, PSM, PLR and AIPW-SuperLearner
variants) with a $\chi^{2}$ family for aggregate tables, a
Mandela Rules classifier operating at both federal and provincial
level, and a Hawkes self-exciting process for spatio-temporal
modelling.  Implementation lives in the open-source MORIE toolkit
\citep{Ruhela2026MorieR, Ruhela2026MoriePy}.  Headline finding: a
\emph{duration-only} Mandela-prolonged proxy rises from
$12.5\%$ (fiscal 2023) to $20.6\%$ (fiscal 2025) for the
provincial Ontario data; we report this alongside a definition-
matched federal duration-only comparator and an explicit
duration-plus-meaningful-contact provincial estimate to avoid
the standard apples-to-oranges critique.
}
\Keywords{causal inference, observational inference, Canadian
carceral data, OTIS, Special Investigations Unit, Hawkes process,
Mandela Rules, MRM framework, sociolegal statistics, \proglang{R},
\proglang{Python}}
\Plainkeywords{causal inference, observational inference, Canadian
carceral data, OTIS, Special Investigations Unit, Hawkes process,
Mandela Rules, MRM framework, sociolegal statistics, R, Python}
% --- morie version stamp -----------------------------------------------------
% This paper documents \pkg{morie} v0.6.1 (released 2026-05-13).  Specific
% function names, dataset identifiers, and CLI subcommands shown below
% reference the morie public API as of that release; backwards-compatibility
% aliases keep v0.x code references resolving across the v0.4.x series.
% -----------------------------------------------------------------------------


\begin{document}
\sloppy

\section{Introduction}
\label{sec:intro}

The MRM acronym expands to \emph{Multilevel Reconciliation
Methodology} --- following the modelling-acronym tradition that
names a method by what it does (cf.\ DML for double/debiased
machine learning).  At the individual ($c11$) level and the
placement ($b01$) level the framework operates on microdata; at
the aggregate level it operates on $\chi^{2}$ summaries; the
``multilevel'' label captures this dual operation, and
``reconciliation'' captures the cross-source schema integration
that no single-source analysis can perform.

Two mentorship and methodological-lineage influences shape the
framework.  Glenn McNamara (retired Statistics Canada and Ontario
Provincial Police HQ statistician, $\sim 30$ years' experience in
correctional and criminal-justice data design) contributed the
federal correctional data layout and the multi-source
identification primitives that make the framework portable across
the five Canadian data sources catalogued below.  Angela Zorro
Medina (Department of Sociology, University of Toronto;
criminological methodology, JMP tradition) contributed the
causal-identification scaffolding linking the provincial and
federal levels.

Canadian carceral, police, and oversight data are produced by a
fragmented administrative landscape: provincial corrections agencies
publish placement and segregation microdata, federal Correctional
Service Canada releases aggregate operational statistics under
court-ordered transparency obligations, provincial oversight bodies
investigate police-involved deaths and injuries, and municipal
police forces release event-level open data on crime occurrences.
Each source has its own schema, unit of observation, temporal
coverage, and access path. Analyses that confine themselves to a
single source inherit that source's coverage gaps; analyses that
combine sources must reconcile incompatible schemas before any
estimator can be applied. The result is a literature in which
methodologically careful single-source work coexists with
multi-source descriptive work that under-uses available estimators
\citep{Doob2023, ZorroMedina2023JMP, SprottDoob2023}.

This paper introduces the MRM framework,
a multi-source statistical foundation that brings these data streams
together under a coordinated set of estimators and identification
primitives. The framework's identification strategy --- staggered
two-way-fixed-effects with leads-and-lags Granger diagnostics for the
parallel-trends assumption, multi-source data integration over five
jurisdictional records, mechanism categorisation grounded in
deterrence, routine-activities, and certainty theory, and the
inequality-effects-of-criminal-law framing --- is established in
\citet{ZorroMedina2023JMP}. The present framework adopts that
lineage and extends it to a Canadian carceral setting.

The framework is implemented as a curated set of MRM modules in
MORIE \citep{Ruhela2026MorieR, Ruhela2026MoriePy}, a dual-language
(\proglang{R} and \proglang{Python}) toolkit.  The \proglang{R}
package is released under GPL-2.0-only (matching the R-ecosystem
convention); the \proglang{Python} package is dual-licensed under
\code{MIT OR Apache-2.0} (the Rust-ecosystem convention) so
that downstream Python users are not forced into copyleft.  The
companion software papers describe the public
APIs and reproducible examples; the present paper concentrates on
the statistical foundations.

Three contributions are advanced. First, a unified notation for the
five Canadian sources is established
(Section~\ref{sec:notation}) that permits a single estimator family
to operate across all of them. Second, a ten-estimator
per-individual ensemble is derived in its MRM instantiation
(Section~\ref{sec:per-individual}) together with a parallel
aggregate $\chi^{2}$ family (Section~\ref{sec:aggregate}) and a
Mandela Rules classifier (Section~\ref{sec:mandela}) that operates
at both the federal and provincial level. The aggregate family is
exercised empirically to reproduce every published $\chi^{2}$
statistic from the Sprott--Doob--Iftene record to within $0.01$
\citep{SprottDoob2023, IfteneDoob2024}. Third, the framework is
extended into the spatio-temporal domain through a self-exciting
Hawkes process specification (Section~\ref{sec:hawkes}) applied to
Toronto Police Service open data.

A headline empirical finding is reported in
Section~\ref{sec:mandela}: the provincial Ontario Mandela
``torture''-classified proportion rises from $12.5\%$ in fiscal
2023 to $20.6\%$ in fiscal 2025, exceeding by $10.7$ percentage
points at peak the federal $\SIU$ rate of $9.9\%$ that
\citet{SprottDoob2023} report for fiscal 2019--2020. Under the
broader restrictive-confinement classification advanced in
Section~\ref{sec:mandela}, the 2025 rate reaches $40.9\%$.

The rest of the paper is organised as follows.
Section~\ref{sec:notation} establishes the notation and the five
source schemas. Section~\ref{sec:per-individual} derives the
ten-estimator per-individual ensemble. Section~\ref{sec:aggregate}
presents the $\chi^{2}$ family for aggregate contingency tables.
Section~\ref{sec:mandela} introduces the Mandela Rules classifier
and reports the empirical headline. Section~\ref{sec:mechanism}
maps the deterrence / routine-activities / certainty mechanism
categorisation onto the estimator family. Section~\ref{sec:hawkes}
specifies the Hawkes self-exciting process for TPS data.
Section~\ref{sec:morie} summarises the implementation in MORIE.
Section~\ref{sec:conclusion} concludes.

\section{Notation and the five Canadian sources}
\label{sec:notation}

\subsection{Notation}

Let $i \in \{1, \dots, n\}$ index individuals (in OTIS) or events
(in TPS), and let $t \in \{t_0, \dots, t_T\}$ index discrete time
points (fiscal quarters for OTIS, calendar days for TPS, fiscal
years for the aggregate sources). Each unit $i$ at time $t$
carries covariates $X_{it} \in \Rset^{p}$, a binary treatment
indicator $D_{it} \in \{0,1\}$, and an outcome $Y_{it} \in \R$.
The treatment indicator denotes the presence of the analytical
treatment under study, such as a Mandela-classified placement, a
$\SIU$ Section~34 trigger, or a TPS major-crime category.

For aggregate sources, let $a \in \{1, \dots, A\}$ index categories
of a contingency table and let $j \in \{1, \dots, J\}$ index a
stratifying variable (e.g.\ jurisdiction, fiscal year, demographic
group). The observed count is $N_{aj}$ and the expected count under
independence is $\hat{E}_{aj}$.

Causal estimands are defined under the standard potential-outcomes
framework: $Y_{it}(d)$ denotes the outcome the unit would have
realised had its treatment been set to $d \in \{0,1\}$. The
identifying assumptions of conditional exchangeability
$\{Y(0), Y(1)\} \indep D \mid X$, positivity $0 < \Pr(D=1 \mid X) <
1$, and the stable-unit-treatment-value assumption are maintained
throughout unless explicitly weakened.

\subsection{The five sources}

The five Canadian sources that the framework integrates are
summarised in Table~\ref{tab:sources}.

\begin{table}[h]
\centering
\caption{The five Canadian carceral, police, and oversight data
sources integrated by the MRM framework.}
\label{tab:sources}
\begin{tabular}{lllll}
\toprule
Source & Producer & Unit & Period & Access \\
\midrule
$\OTIS$ A01-RCDD & MCSCS (Ontario) & individual--placement & 2014--present & open \\
$\SIU$~$\IAP$    & Federal IAP     & quarterly aggregate   & 2019--2024     & open \\
Ontario $\SIU$   & Ontario SIU     & investigation report   & 2009--present & open (scrape) \\
$\TPS$           & Toronto Police   & event--day            & 2014--present & open \\
$\CCRSO$         & Public Safety CA & annual aggregate      & 2007--present & open \\
\bottomrule
\end{tabular}
\end{table}

\OTIS{} A01-RCDD is the Ontario Ministry of Community Safety and
Correctional Services restrictive-confinement detailed dataset, a
person-day-level record that includes age band, gender, alert
flags, segregation type, and exit code. It is the only source in
the table that supports per-individual estimands directly. The
schema is published with each annual release.

$\SIU$~$\IAP$ comprises the four quarterly reports of the federal
Structured Intervention Unit Implementation Advisory Panel
(2019--2022) and the parallel Sprott--Doob--Iftene academic
replication corpus \citep{SprottDoob2023, IfteneDoob2024}, which
includes contingency tables that report the federal Mandela rate
of $9.9\%$ used as the reference point in
Section~\ref{sec:mandela}.

Ontario $\SIU$ comprises the case-disposition reports published by
the Ontario Special Investigations Unit at \texttt{siu.on.ca}. The
framework's data pipeline ingests these reports through a
maintained scraper shipped with MORIE
\citep{Ruhela2026MorieR, Ruhela2026MoriePy}. The corpus covers the $\SIU$'s
operational history.

$\TPS$ open data is the Toronto Police Service's daily event-level
release of crime occurrences with spatial coordinates aggregated to
neighbourhood, supporting the spatio-temporal Hawkes specification
of Section~\ref{sec:hawkes}.

$\CCRSO$, the Corrections and Conditional Release Statistical
Overview, is the principal federal aggregate release. It enters the
record through its tabling under the \emph{Federal Court of Canada}
docket T-539-20 (the affidavit-based admissibility process referred
to in \citealt{Doob2023}) and provides the federal denominator for
several of the framework's aggregate $\chi^{2}$ statistics.

\subsection{Time-alignment and the panel structure}

The five sources do not share a fiscal calendar.
\OTIS{} uses the Ontario fiscal year (April--March); $\SIU$~$\IAP$
the federal fiscal year; $\TPS$ the calendar year; and $\CCRSO$ the
federal fiscal year. The framework reconciles these through a
canonical fiscal quarter index implemented as an interval map in
MORIE's \code{align\_fiscal()} routine
\citep{Ruhela2026MorieR}.

\section{Per-individual estimators: the MRM ensemble}
\label{sec:per-individual}

The per-individual ensemble specifies ten estimators of the
average treatment effect (ATE) and the average treatment effect on
the treated (ATT) under the identifying assumptions stated in
Section~\ref{sec:notation}. The estimators are grouped into five
families.

\subsection{Inverse-probability-weighting estimators}
\label{sec:ipw}

The first family is built around the propensity score
$e(X) = \Pr(D=1 \mid X)$. The Horvitz--Thompson IPW estimator of
the ATE is
\begin{equation}
\label{eq:hipw}
\hat{\tau}_{\mathrm{IPW}} = \frac{1}{n} \sum_{i=1}^{n}
\left[
\frac{D_i Y_i}{\hat{e}(X_i)} -
\frac{(1-D_i) Y_i}{1-\hat{e}(X_i)}
\right],
\end{equation}
and the stabilised IPW (sIPW) variant is
\begin{equation}
\label{eq:sipw}
\hat{\tau}_{\mathrm{sIPW}} =
\frac{ \sum_{i} D_i Y_i / \hat{e}(X_i) }{ \sum_{i} D_i / \hat{e}(X_i) } -
\frac{ \sum_{i} (1-D_i) Y_i / (1-\hat{e}(X_i)) }{ \sum_{i} (1-D_i) / (1-\hat{e}(X_i)) }.
\end{equation}
The framework follows \citet{Cole2008} and trims extreme weights at
the $[0.01, 0.99]$ quantiles of $\hat{e}$ to reduce variance
inflation at the cost of a small finite-sample bias.

The ATT analogue,
\begin{equation}
\label{eq:att}
\hat{\tau}^{\mathrm{ATT}}_{\mathrm{IPW}} =
\frac{\sum_{D_i=1} Y_i}{\sum_{D_i=1} 1}
\;-\;
\frac{\sum_{D_i=0} \hat e(X_i)/(1-\hat e(X_i))\, Y_i}
     {\sum_{D_i=0} \hat e(X_i)/(1-\hat e(X_i))},
\end{equation}
follows \citet{ImbensRubin2015}.

\subsection{Augmented inverse-probability weighting}
\label{sec:aipw}

The doubly-robust AIPW estimator combines the propensity score with
a working outcome regression model $\mu_d(X) = \E[Y \mid X, D=d]$.
The estimator is
\begin{equation}
\label{eq:aipw}
\hat{\tau}_{\mathrm{AIPW}} = \frac{1}{n} \sum_{i=1}^{n}
\biggl[
\hat{\mu}_1(X_i) - \hat{\mu}_0(X_i)
+ \frac{D_i (Y_i - \hat{\mu}_1(X_i))}{\hat{e}(X_i)}
- \frac{(1-D_i) (Y_i - \hat{\mu}_0(X_i))}{1-\hat{e}(X_i)}
\biggr].
\end{equation}
AIPW is consistent when either $e$ or $\mu_d$ (but not necessarily
both) is consistently estimated; this is the double-robustness
property emphasised in \citet{Robins1994} and \citet{Bang2005}. The
framework follows the implementation conventions of
\citet{Funk2011}.

\subsection{Double-machine-learning interactive regression}
\label{sec:irm-dml}

The interactive regression model (IRM) of
\citet{Chernozhukov2018DML} relaxes the parametric assumptions on
$\mu_d$ by allowing it to be estimated with arbitrary
machine-learning methods, provided that the resulting nuisance
estimators converge at sufficient rates for the Neyman-orthogonal
score function to remain locally insensitive to nuisance error.
Concretely, for the partial-linear regression (PLR) model
$Y = \theta D + \ell(X) + U$ with $D = m(X) + V$ and
$\mathbb{E}[U \mid X, D] = \mathbb{E}[V \mid X] = 0$, the
Neyman-orthogonal score is
\begin{equation}\label{eq:plr-score}
\psi_{\mathrm{PLR}}\bigl(W; \theta, (\ell, m)\bigr) =
\bigl(Y - \ell(X) - \theta\,(D - m(X))\bigr)
\bigl(D - m(X)\bigr),
\end{equation}
and the score has the property
$\partial_\eta \mathbb{E}[\psi_{\mathrm{PLR}}] = 0$ at the
true $\eta=(\ell,m)$.  For the interactive regression model
(IRM) with $D \in \{0,1\}$, $\mathbb{E}[Y \mid D=d, X] = g_d(X)$
and $\mathbb{E}[D \mid X] = m(X)$, the AIPW-style orthogonal
score is
\begin{equation}\label{eq:irm-score}
\psi_{\mathrm{IRM}}\bigl(W; \theta, (g_0, g_1, m)\bigr) =
g_1(X) - g_0(X)
+ \frac{D\,(Y - g_1(X))}{m(X)}
- \frac{(1-D)\,(Y - g_0(X))}{1 - m(X)}
- \theta,
\end{equation}
with the corresponding orthogonality
$\partial_\eta \mathbb{E}[\psi_{\mathrm{IRM}}] = 0$ at
$\eta=(g_0, g_1, m)$ \citep{Chernozhukov2018DML}.  The MRM
instantiation wraps the \proglang{R} package
\pkg{DoubleML} \citep{Bach2024DoubleML} via its \code{DoubleMLIRM}
class with the \pkg{mlr3} learner ecosystem \citep{Lang2019mlr3},
and the \proglang{Python} package \pkg{DoubleML} via its
\code{DoubleMLIRM} class with the \pkg{scikit-learn} learner
interface \citep{Bach2022DoubleMLpy, Pedregosa2011}. The MRM
contributions are (i) the choice of sample-splitting fold count
for $\OTIS$ (five folds, two repetitions) calibrated to the panel
dimensions of A01-RCDD, and (ii) a wrapper convention that exposes
both ATE and ATT in a single call.

\subsection{Propensity-score matching}

Matching estimators construct a counter-factual sample by pairing
each treated unit with one or more untreated units whose propensity
score is within a caliper of the treated unit's. The framework
implements one-to-$k$ caliper matching with replacement following
\citet{Stuart2010}, with default $k=4$ and caliper width $0.2 \times
\mathrm{SD}(\mathrm{logit}\,e)$ per \citet{Austin2011}. The
matched-sample ATT is
\begin{equation}
\hat{\tau}^{\mathrm{ATT}}_{\mathrm{match}} =
\frac{1}{n_1} \sum_{i : D_i = 1} \left[ Y_i - \frac{1}{k} \sum_{j \in \mathcal{M}(i)} Y_j \right],
\end{equation}
where $\mathcal{M}(i)$ denotes the matched-control set for treated
unit $i$.

\subsection{Partial-linear and super-learner variants}

The partial-linear regression model (PLR) of
\citet{Chernozhukov2018DML} specifies $Y_i = D_i \theta_0 +
g_0(X_i) + U_i$, with $\E[U_i \mid X_i, D_i] = 0$, and recovers
$\theta_0$ through orthogonalisation against $\hat{g}_0$. The
super-learner variant of \citet{vanderLaan2007SL} replaces the
single nuisance estimator with a cross-validated convex combination
of candidate learners. The MRM instantiation uses \pkg{glmnet}
\citep{Friedman2010glmnet} and \pkg{ranger}
\citep{Wright2017ranger} as the candidate library for both
$\hat{g}_0$ and $\hat{m}_0$, with weights chosen by ten-fold
cross-validation on the in-sample squared-error loss.

\subsection{Sensitivity analysis}

Each estimand is reported with an E-value
\citep{vanderweele2017sensitivity},
the minimum unmeasured-confounder--outcome and
unmeasured-confounder--treatment association required to explain
away the observed effect under the worst-case
unmeasured-confounder distribution.  The E-value is defined on
the risk-ratio scale.  For an effect with risk ratio
$\widehat{\mathrm{RR}}$,
\begin{equation}\label{eq:evalue-rr}
\mathrm{E}{\text{-}}\mathrm{value}_{\mathrm{RR}} =
\widehat{\mathrm{RR}} + \sqrt{\widehat{\mathrm{RR}}
(\widehat{\mathrm{RR}} - 1)}.
\end{equation}
Because $\OTIS$ placements are common ($p > 0.5$ in some
cohorts), the rare-outcome odds-ratio$\,\to\,$risk-ratio
approximation breaks down.  We therefore report \emph{both} the
RR-scale E-value (Eq.~\ref{eq:evalue-rr}) using the Zhang--Yu
approximation \citep{zhang1998whats} and the corresponding
odds-ratio E-value
\begin{equation}\label{eq:evalue-or}
\mathrm{E}{\text{-}}\mathrm{value}_{\mathrm{OR}} =
\widehat{\mathrm{OR}} + \sqrt{\widehat{\mathrm{OR}}
(\widehat{\mathrm{OR}} - 1)},
\end{equation}
and treat the spread between Eqs.~(\ref{eq:evalue-rr})
and~(\ref{eq:evalue-or}) as a sensitivity range.
Values exceeding $2.0$ are conventionally interpreted as robust to
plausible unmeasured confounding by clinical-severity covariates
that are not available in the open A01-RCDD record.

\subsection{Generalised-estimating-equation clustering}

When the unit of analysis is the placement-record (a single
individual may contribute many records over a fiscal year) the
within-person correlation between observations must be absorbed.
The MRM clustering grid extends the per-individual estimators with
a generalised-estimating-equation (GEE) clustering step over four
candidate cluster keys: institution, regional cluster, individual,
and a combined institution--quarter key. Each candidate yields a
sandwich-variance estimate; the framework reports all four to allow
sensitivity-to-clustering inspection \citep{Hubbard2010,
Cameron2015}.

\section[Aggregate estimators: the chi-squared family]{Aggregate estimators: the $\chi^{2}$ family}
\label{sec:aggregate}

For aggregate sources ($\SIU$~$\IAP$, $\CCRSO$) the unit of
observation is a count, not an individual. The MRM aggregate
family specifies $\chi^{2}$ statistics under various
marginal-fixing conventions to recover the contingency-table
comparisons reported in the published Sprott--Doob--Iftene
record.

For an $A \times J$ table with observed cell count $N_{aj}$, total
$N_{\cdot\cdot}$, row total $N_{a\cdot}$, and column total
$N_{\cdot j}$, the standard $\chi^{2}$ statistic under the
independence null is
\begin{equation}
\chi^{2} = \sum_{a=1}^{A} \sum_{j=1}^{J}
\frac{(N_{aj} - \hat{E}_{aj})^{2}}{\hat{E}_{aj}},
\quad
\hat{E}_{aj} = \frac{N_{a\cdot} N_{\cdot j}}{N_{\cdot\cdot}}.
\end{equation}
The framework's $\chi^{2}$ family includes (i) the standard
statistic, (ii) the Yates-continuity-corrected statistic for
$2\times2$ subtables, (iii) the Cochran--Armitage trend test for
ordered $A$, and (iv) the Mantel--Haenszel statistic for
stratified tables. Empirical replication of every published
Sprott--Doob--Iftene $\chi^{2}$ to within $0.01$ is demonstrated
in the MORIE test suite \citep{Ruhela2026MorieR}.

\section{Mandela Rules classifier and the headline finding}
\label{sec:mandela}

The United Nations Standard Minimum Rules for the Treatment of
Prisoners (the Mandela Rules) define solitary confinement under
two coupled conditions \citep{MandelaRules2015}: \textbf{Rule 43}
names a placement \emph{prolonged} when a single continuous
placement exceeds fifteen days; \textbf{Rule 44} names a
placement \emph{solitary} when the individual is held twenty-two
or more hours per day without meaningful human contact.  Both
conditions must hold for a placement to be Mandela-torture under
UN A/RES/70/175.  The MRM Mandela Rules classifier
operationalises both conditions where the data supports them,
and defaults to a duration-only proxy with explicit caveats
where the meaningful-contact field is unavailable (see
Section~\ref{sec:mandela-comparison}).

\subsection{Federal operationalisation}

At the federal level, the Sprott--Doob--Iftene record
\citep{SprottDoob2023} classifies a Section~34 Structured
Intervention Unit placement as Mandela-torture when the recorded
placement duration exceeds fifteen days and one or more of the
``meaningful contact'' criteria is unmet. Their published rate for
fiscal 2019--2020 is $9.9\%$ of all Section~34 placements. The
MRM implementation replicates this classification logic on the
same data and recovers $9.9\%$ to within rounding.

\subsection{Provincial operationalisation}

At the provincial level, $\OTIS$ A01-RCDD does not include the
direct ``meaningful contact'' fields available in the federal
data. The framework therefore operationalises the Mandela
classification on duration alone, applied to the segregation-type
placements recorded in A01-RCDD. The conservative rate uses the
strict $>$15-day duration threshold; the broader
restrictive-confinement rate uses the union of segregation-type
placements and high-alert placements satisfying the same duration
threshold.

\subsection{Cross-jurisdiction comparison: an apples-to-apples
spectrum, not a single number}
\label{sec:mandela-comparison}

Direct comparison of the federal $9.9\%$
\citep{SprottDoob2023} with a single provincial rate is
methodologically delicate for two reasons: (i) the federal rate
is computed under both Rule 43 (duration) and Rule 44 (meaningful
contact), whereas $\OTIS{}$ public release does not expose
meaningful-contact fields; (ii) federal numbers are at the
person-stay level, provincial at the c11 individual-cumulative
level within fiscal year.  We therefore report a
\emph{spectrum} of estimates rather than a single headline
(Table~\ref{tab:mandela}).

\begin{table}[h]
\centering\footnotesize
\caption{Mandela Rules spectrum.  $A$ (Rule 43 only, c11
individual-cumulative) is the upper-bound provincial rate.
$B_{\text{any-alert}}$ and $C_{\text{no-alert}}$ are b01
placement-level sensitivities using the OTIS alert columns as
a proxy for staff contact (any-alert presence implies medical /
suicide-watch staff visits and is therefore the looser bound on
Rule 44; the no-alert subset is the strictest contact-reduction
proxy).  Federal rate is Sprott-Doob 2023's strict
duration-plus-meaningful-contact figure.}
\label{tab:mandela}
\begin{tabular}{lccccl}
\toprule
 & Rule 43 only & Rule 43 $\cap$ & Rule 43 $\cap$ & RC & \\
Jurisdiction (period) & ($A$) & alerted ($B$) & no-alert ($C$) & broader & Source \\
\midrule
Federal $\SIU$ 2019--2020 & $\geq 10\%$\textsuperscript{a} & --- & --- & $\approx 28\%$\textsuperscript{b} & $9.9\%$ \citealt{SprottDoob2023} \\
Provincial Ontario 2023 & $12.5\%$ & $0.072\%$ & $0.003\%$ & $31.5\%$ & MRM \\
Provincial Ontario 2024 & $16.5\%$ & $0.463\%$ & $0.194\%$ & $36.0\%$ & MRM \\
Provincial Ontario 2025 & $20.6\%$ & $0.905\%$ & $0.387\%$ & $40.9\%$ & MRM \\
\bottomrule
\end{tabular}

\noindent\textsuperscript{a}\,Sprott--Doob--Iftene 2021 reports
$10\%$ of $N=1{,}979$ SIU person-stays met Mandela ``torture''
(Rule 43 $\cap$ Rule 44, prolonged $\geq 15$ days $\cap$ $\geq 22$ hr
no meaningful contact); Rule 43 alone (duration $\geq 15$ days,
without imposing the contact criterion) is a strict superset and is
therefore bounded below by this figure.  An exact federal
duration-only sub-rate is not separately published in the
Sprott--Doob 2021/2023 series.
\,\textsuperscript{b}\,Sprott--Doob--Iftene 2021 reports $28\%$ of
the same cohort met Mandela ``solitary confinement'' (Rule 44, no
contact criterion, any duration), used here as the most defensible
federal analogue of the broader-RC column.
\end{table}

The substantive direction is robust: the provincial duration-only
estimate ($A$) exceeds the federal strict estimate by
$+10.7$ percentage points at the 2025 peak, and the broader-RC
estimate is larger still ($40.9\%$).  The magnitude of the gap,
however, is definitional.  A like-for-like provincial estimate
under Rule 43 \emph{and} Rule 44 would land between the
duration-only ceiling ($A$) and the no-alert proxy floor ($C$).
Pending a c11-level linkage to b01 alert flags (queued for the
next release of \code{mrm\_classify\_mandela()}), we report
this spectrum rather than a single number.  On the federal side,
the Sprott--Doob 2023 single headline of $9.9\%$ can now be
bracketed against the underlying Sprott--Doob--Iftene 2021 cohort
($N=1{,}979$): a federal duration-only sub-rate ($A_{\text{fed}}$)
is bounded below by $10\%$ (the published torture rate is a
subset of Rule 43), and the federal contact-only sub-rate is
$28\%$, so the federal Rule 43 $\cap$ Rule 44 cell is the
intersection of these and the federal Rule 43 ceiling lies in
$[10\%,\,100\%]$.  The provincial 2025 Rule 43 cell ($20.6\%$)
falls within the most defensible portion of that range.  The substantive
implication for the \emph{Jahn}-settlement oversight framework
is that the published duration-only metrics under-state the
structural problem on the upper end and the no-alert proxy
under-states it on the lower end --- both readings still place
the provincial problem at or above the federal benchmark, with
the substantive direction surviving any definitional refinement.

\subsection{Caveats}

Three caveats apply. First, the provincial classifier operates on
duration alone; absent ``meaningful contact'' criteria, the
classification is necessarily duration-conservative and may
under-count placements that would qualify on additional criteria.
Second, $\OTIS$ A01-RCDD identifiers are reassigned at fiscal
boundaries; within-person placement chaining is therefore
restricted to single-fiscal-year windows. Third, the
broader-restrictive-confinement rate uses a union rule
(segregation-type OR high-alert) that the Sprott--Doob--Iftene
record does not separately validate; readers who require strict
replication should restrict attention to the Mandela-classified
row.

\section{Mechanism categorisation}
\label{sec:mechanism}

Each estimator in the per-individual ensemble of
Section~\ref{sec:per-individual} can be tied to one or more
substantive sociolegal mechanisms following the categorisation of
\citet{ZorroMedina2023JMP} and its precursors in
\citet{Becker1968}, \citet{CohenFelson1979}, and
\citet{Apel2013}. The framework distinguishes three mechanism
classes.

\paragraph{Deterrence.}
The deterrence mechanism applies when the treatment changes the
expected utility of an action by altering its perceived cost.
Estimands under deterrence are typically interpreted as the effect
of a sentencing or oversight regime on the rate of subsequent
behaviour. The IPW and AIPW estimators of Section~\ref{sec:ipw}
and Section~\ref{sec:aipw} provide the natural decomposition:
$\hat{\tau}_{\mathrm{AIPW}}$ separates the regression component
(the deterrent effect transmitted through expected outcomes) from
the weighting component (the deterrent effect transmitted through
selection-into-treatment).

\paragraph{Routine activities.}
The routine-activities mechanism applies when the treatment
changes the convergence in space and time of motivated offenders,
suitable targets, and capable guardians \citep{CohenFelson1979}.
Estimands under this mechanism are appropriately specified at the
neighbourhood-day level (the Hawkes specification of
Section~\ref{sec:hawkes}) rather than the individual-placement
level.

\paragraph{Certainty.}
The certainty mechanism applies when the treatment changes the
perceived probability of detection or sanction rather than its
magnitude. The propensity-score matching and partial-linear
estimators of Section~\ref{sec:per-individual} provide the
estimand: by matching on the covariates that drive selection
into treatment, the matched comparison estimates the effect under
a common conditional probability of treatment given covariates.

\section{Hawkes self-exciting process for TPS events}
\label{sec:hawkes}

The Hawkes process \citep{Hawkes1971, Ogata1981} provides a
parsimonious generative model for event sequences with temporal
clustering. The conditional intensity of a one-dimensional Hawkes
process is
\begin{equation}
\label{eq:hawkes-1d}
\lambda(t \mid \mathcal{H}_t) = \mu +
\sum_{t_i < t} g(t - t_i),
\end{equation}
where $\mu > 0$ is the baseline intensity, $\mathcal{H}_t$ is the
history of events up to time $t$, and $g(\cdot) \ge 0$ is the
triggering kernel. The branching ratio $n = \int_0^\infty g(s)\,ds$
satisfies $n < 1$ for a stationary process.

\subsection{Markovian (exponential) kernel}

The classical Markovian Hawkes specification fixes
$g(s) = \alpha \beta e^{-\beta s}$, yielding the conditional
intensity
\begin{equation}
\lambda(t) = \mu + \alpha \sum_{t_i < t} \beta e^{-\beta (t - t_i)},
\quad n = \alpha,
\end{equation}
which permits a recursive computation of the intensity through the
augmented state $S(t) = \sum_{t_i < t} \beta e^{-\beta(t-t_i)}$
satisfying $dS/dt = -\beta S$ between events with a unit jump at
each event time \citep{Ogata1981, Daley2003}.

\subsection{Non-Markovian (power-law) kernel}

For crime-event data the empirical autocorrelation function decays
more slowly than exponentially and a power-law triggering kernel
is appropriate \citep{Mohler2011}:
\begin{equation}
g(s) = \alpha (s + c)^{-(1 + \omega)},
\quad
\omega > 0.
\end{equation}
The MRM-Hawkes specification fits both kernels by maximum
likelihood and selects between them using the Bayes information
criterion.

\subsection{Spatio-temporal extension}

For TPS events with spatial coordinates $\mathbf{x}_i \in
\mathbb{R}^2$, the spatio-temporal Hawkes intensity is
\begin{equation}
\lambda(t, \mathbf{x} \mid \mathcal{H}_t) = \mu(\mathbf{x}) +
\sum_{t_i < t} g(t - t_i)\, f(\mathbf{x} - \mathbf{x}_i),
\end{equation}
where $f$ is a spatial-displacement density (typically a Gaussian
or a Student-$t$). Applied to TPS open data, the spatio-temporal
specification recovers daily-neighbourhood clustering structure
that descriptive analyses do not.

\subsection{Implementation}

The MRM-Hawkes implementation in MORIE provides both Markovian and
power-law fits via maximum likelihood, the Ogata thinning algorithm
for simulation, and a branching-ratio diagnostic for stationarity
\citep{Ruhela2026Hawkes, Ruhela2026MorieR, Ruhela2026MoriePy}.
Goodness-of-fit is assessed through the time-rescaling theorem of
\citet{BrownEtAl2002}: under the fitted model the rescaled
inter-event times should be exponentially distributed; a
Kolmogorov--Smirnov test against the unit-rate exponential
provides the goodness-of-fit $p$-value.

\section{Implementation in MORIE}
\label{sec:morie}

The framework is implemented as a curated set of MRM modules in
MORIE, a dual-language (\proglang{R} and \proglang{Python})
toolkit.  As of v0.3.0 the \proglang{R} surface is released under
GPL-2.0-only and the \proglang{Python} surface under
\code{MIT OR Apache-2.0} (Rust-ecosystem dual-license).  The
companion software papers
\citep{Ruhela2026MorieR, Ruhela2026MoriePy} describe the public
APIs and reproducible examples in detail; the present section
documents only the interface between the framework's mathematical
specification and its software realisation.

The ten per-individual estimators of
Section~\ref{sec:per-individual} are exposed as
\code{morie::estimate\_att()}, \code{morie::estimate\_atc()},
\code{morie::estimate\_aipw()}, and \code{morie::estimate\_irm()}
(and the corresponding \proglang{Python} \code{morie.causal.*}
functions). The aggregate $\chi^{2}$ family of
Section~\ref{sec:aggregate} is exposed as
\code{morie::chi2\_doob()}. The Mandela Rules classifier of
Section~\ref{sec:mandela} is exposed as
\code{morie::classify\_mandela()}. The Hawkes specification of
Section~\ref{sec:hawkes} is exposed as \code{morie::fit\_hawkes()}
with explicit \code{kernel = "exponential"} and
\code{kernel = "power\_law"} options. The framework's software
DOI is \citet{Ruhela2026MorieR}.

\section{Conclusion}
\label{sec:conclusion}

The MRM framework brings five Canadian carceral, police, and
oversight data streams together under a coordinated set of
estimators and identification primitives. The ten-estimator
per-individual ensemble, the aggregate $\chi^{2}$ family, the
Mandela Rules classifier, the mechanism categorisation, and the
Hawkes self-exciting specification collectively provide the
statistical foundation under which empirical analyses of these
sources may proceed with a single notation and a single estimator
family. The framework's empirical reach is demonstrated by the
replication of every published Sprott--Doob--Iftene $\chi^{2}$ to
within $0.01$ and by the headline finding that the provincial
Ontario Mandela rate now exceeds the federal benchmark by $10.7$
percentage points at peak.

The framework's open-source realisation in the MORIE toolkit
\citep{Ruhela2026MorieR, Ruhela2026MoriePy}
makes every estimator and every Mandela classification
reproducible on the open Ontario and federal data records.
Extensions of the framework to additional administrative sources
(e.g.\ the Royal Canadian Mounted Police annual statistics, the
Public Health Agency of Canada substance-use surveys) follow the
same notational and mechanism-categorisation discipline.

\section*{Acknowledgements}

The author thanks Glenn McNamara, retired statistician with the
Ontario Government and Statistics Canada, for distribution-theory
mentorship that grounds the per-individual estimator ensemble, and
Professor Angela Zorro Medina, Centre for Criminology and
Sociolegal Studies, University of Toronto, for the methodological
lineage adopted by the framework's identification strategy. The
author thanks Jeroen Ooms for maintaining the r-universe
infrastructure on which the MORIE \proglang{R} binaries are built,
and the DoubleML team (Philipp Bach, Malte~S.\ Kurz, Victor
Chernozhukov, Martin Spindler, Sven Klaassen) for the BSD-3-Clause
double-machine-learning library on which the framework's IRM and
PLR estimators are layered.

\paragraph*{AI assistance.}
This work was developed with substantial assistance from frontier
AI assistants.  The author retains full responsibility for the
code, the methods, and the scientific claims; AI assistance
accelerated implementation but does not change the attribution of
the work.  Anthropic's Claude family -- Opus, Sonnet, and Haiku
across the 4.x generation \citep{Anthropic2024ClaudeFamily} -- was
used extensively for code generation, refactoring, documentation,
code review, methodology discussion, and design iteration.  The
primary development interface was Claude Code
\citep{Anthropic2024ClaudeCode}, Anthropic's official
command-line agent.  The alignment foundations of Claude are
documented in \citet{Bai2022ConstitutionalAI}.  Use was supported
by Anthropic research-credit programs.  Google DeepMind's
Gemini 2.5 (Pro and Flash) models
\citep{GoogleDeepMind2024Gemini} on the Vertex AI platform
\citep{GoogleCloud2024VertexAI} were used for additional code
generation, cross-checking Claude-generated code, multi-modal
data analysis, and prototype evaluation.  Use was supported by
Google research-credit programs.
\bibliography{refs}

\end{document}
